\begin{abstract}
% Numbers: tables/tab_flux_main.tex and tables/tab_video_main.tex (every row run in the same setup);
% research_v2/runs/drawbench_eval/eval_dct_pcgs.txt (KalSea, delta 0.37 / 0.6); PCGS alone: paired vs. the predictor
% (speckcast_<d>), figures/src/data/flux_drawbench_stats.json (eval_sweep_vs_kalsea.json per-image metrics, eval2 bootstrap);
% gains on reuse / linear extrapolation: research_v2/ablations/eval/eval_all.json.
Diffusion transformers spend most of their sampling time in the block stack, which runs at each of tens of sampling
steps. Training-free caches skip it on some steps: a gate decides \emph{when} to compute, and a filler decides
\emph{what} to write on the skipped steps, typically the last block-stack residual or a forecast of it. Both fillers
are open loop: when the next computed step reveals how wrong the filler was, its error has already been integrated
into the latent and is never corrected. We close this loop with predictor-corrector gap smoothing (\pcgs). Behind the
spectrally aligned gate of \seacache, a damped-trend Kalman filter forecasts the block-stack residual on skipped
steps. When the next step is computed, a Rauch-Tung-Striebel smoother turns the filter's innovation into corrected
residuals for the whole skipped gap, which we re-integrate through the output layers only before sampling continues
from the corrected latent. \pcgs requires no additional transformer evaluations. On \flux with held-out DrawBench prompts, with every baseline run in the same setup, our cache reaches a PSNR
of 27.72\,dB with respect to the 50-step sampler at 18 of 50 transformer evaluations, against 25.33\,dB for
\seacache and 25.48\,dB for Spectrum, and at 21 evaluations it exceeds every baseline that uses more, including SADA.
In paired comparisons, \pcgs alone contributes 0.81\,dB (95\% CI [0.63, 1.00]) and 0.61\,dB [0.48, 0.76], and it
also improves plain reuse and linear extrapolation.
% Video: figures/src/data/video_measured.json (same-setup rows of tables/tab_video_main.tex).
On HunyuanVideo and Wan2.1 it has the best PSNR, SSIM and LPIPS of all methods at both budgets (1.9--3.2 and
0.3--2.3\,dB above \seacache), with VBench within 0.2 points of the uncached model.
\end{abstract}
